%% ====================================================================
%% CIVIC-SAFE -- Supplementary Material
%%
%% Companion to paper/civic_safe_ieee.tex. Section numbers are prefixed S
%% and cross-references to the main text are given by name, not by number,
%% because the two documents are compiled separately.
%%
%% Compile:  pdflatex civic_safe_supplementary -> bibtex -> pdflatex x2
%%
%% Deliberately pure ASCII. Accented names use LaTeX escapes (Cand\`es,
%% P\'olya) so it compiles without inputenc on any distribution.
%% ====================================================================
\documentclass[journal,10pt,onecolumn]{IEEEtran}

\usepackage{amsmath,amssymb,amsfonts}
\usepackage{booktabs}
\usepackage{graphicx}
\usepackage{cite}
\usepackage{url}
\usepackage[hidelinks]{hyperref}

\graphicspath{{../outputs/figures/}{figures/}{./}}

\newtheorem{lemma}{Lemma}
\newtheorem{proposition}{Proposition}

% Number everything with an S prefix so nothing collides with the main text.
\renewcommand{\thesection}{S\arabic{section}}
\renewcommand{\theequation}{S\arabic{equation}}
\renewcommand{\thetable}{S\arabic{table}}
\renewcommand{\thefigure}{S\arabic{figure}}
\renewcommand{\thelemma}{S\arabic{lemma}}
\renewcommand{\theproposition}{S\arabic{proposition}}

\extrafloats{32}

\begin{document}

\title{Supplementary Material for\\
``Feedback-Corrected Conformal Prediction for Spatiotemporal Crime
Forecasting''}

\author{Himanshu~Bairwa%
\thanks{H. Bairwa is with the Department of Computer Science. Corresponding
author e-mail: bairwahimanshu29@gmail.com.}%
}

\maketitle

\begin{abstract}
This document supplies the derivations, protocols and audits that the main
manuscript states without proof or in summary. Section~\ref{sec:s1} gives the
full spectral argument for the log-linearized Jacobian, a complete proof of the
attribution corollary including sign preservation, a treatment of why integer
conformity scores make coverage conservative, and the sensitivity bound for a
mismeasured feedback gain. Section~\ref{sec:s2} tabulates the full
$\Gamma$-sensitivity frontier. Section~\ref{sec:s3} specifies every baseline
architecture and the training protocol, read from the released code rather than
recalled. Section~\ref{sec:s4} audits the demographic covariates and the two
spatial graphs, including a connectivity defect in one city that motivates the
dual-graph design. Section~\ref{sec:s5} reports the diagnostics we can compute
from the persisted artifacts, and states precisely which requested diagnostic we
cannot, and why.
\end{abstract}

% =====================================================================
\section{Extended Derivations}
\label{sec:s1}
% =====================================================================

Throughout, $s \in \{1,\dots,S\}$ indexes spatial units, $\lambda_s > 0$ is the
unobserved latent incidence rate, $\mu_s > 0$ the recorded-rate estimate,
$M = S^{-1}\sum_s \mu_s$ the panel mean, $\kappa \in [0,1)$ the closed-loop
feedback gain, $\varphi$ the allocation policy, and $g$ the recording response.
Notation follows Table~I of the main text exactly. Section~\ref{sec:s1-jacobian}
proves the spectral claim behind the stability threshold,
Section~\ref{sec:s1-attribution} the attribution corollary and its sign
preservation, Section~\ref{sec:s1-lattice} the lattice effect that makes discrete
coverage conservative, and Section~\ref{sec:s1-sensitivity} the error bound for a
mismeasured gain.

\subsection{The log-linearized Jacobian and exact spectral radius}
\label{sec:s1-jacobian}

The main text asserts that the fixed-point iteration contracts precisely when
$\kappa < 1$, with spectral radius exactly $\kappa$ irrespective of how
heterogeneous the cells are. That claim is proved here, because the ``exactly''
is doing real work: a bound of the form $\rho \le \kappa$ would leave open the
possibility that some configurations contract faster and others not at all.

Under the constant-elasticity parameterization the iteration map
$T : \mathbb{R}^S_{>0} \to \mathbb{R}^S_{>0}$ is
\begin{equation}
\label{eq:s-map}
T(\mu)_s \;=\; \lambda_s \left(\frac{\mu_s}{\bar{\mu}}\right)^{\!\kappa},
\qquad \bar{\mu} = S^{-1}\textstyle\sum_j \mu_j .
\end{equation}
Work in logarithmic coordinates $u_s = \log \mu_s$, which is natural here because
the map is a power law and because positivity is then automatic. Writing
$\bar{\mu}(u) = S^{-1}\sum_j e^{u_j}$,
\begin{equation}
\label{eq:s-logmap}
\widetilde{T}(u)_s = \log \lambda_s + \kappa\left(u_s - \log \bar{\mu}(u)\right).
\end{equation}

\begin{lemma}[Exact spectral radius]
\label{lem:spectral}
The Jacobian of \eqref{eq:s-logmap} at any $u$ is
\begin{equation}
\label{eq:s-jacobian}
J \;=\; \kappa\left(I - \mathbf{1}w^{\!\top}\right),
\qquad
w_j \;=\; \frac{e^{u_j}}{\textstyle\sum_k e^{u_k}} \;=\; \frac{\mu_j}{S\,M},
\end{equation}
where $\mathbf{1} \in \mathbb{R}^S$ is the all-ones vector and $w$ lies in the
probability simplex. Its spectrum is $\{\kappa$ with multiplicity $S-1$, and
$0\}$, so $\rho(J) = \kappa$ for every $u$ and every distribution of $\mu$.
\end{lemma}

\begin{IEEEproof}
Differentiating \eqref{eq:s-logmap},
\begin{equation}
\frac{\partial \widetilde{T}(u)_s}{\partial u_j}
= \kappa\left(\delta_{sj} - \frac{\partial \log \bar{\mu}}{\partial u_j}\right)
= \kappa\left(\delta_{sj} - w_j\right),
\end{equation}
since $\partial \bar{\mu}/\partial u_j = S^{-1}e^{u_j}$ and hence
$\partial \log\bar{\mu}/\partial u_j = e^{u_j}/\sum_k e^{u_k} = w_j$. In matrix
form this is \eqref{eq:s-jacobian}, and $w^{\!\top}\mathbf{1} = 1$ by
construction.

For the spectrum, note $P = \mathbf{1}w^{\!\top}$ is rank one, so it has
eigenvalue $w^{\!\top}\mathbf{1} = 1$ with eigenvector $\mathbf{1}$, and
eigenvalue $0$ with multiplicity $S-1$ on the hyperplane
$\{v : w^{\!\top}v = 0\}$. Indeed $P\mathbf{1} = \mathbf{1}(w^{\!\top}\mathbf{1})
= \mathbf{1}$ and $Pv = \mathbf{1}(w^{\!\top}v) = 0$ for such $v$. Therefore
$I - P$ has eigenvalue $0$ on $\mathrm{span}\{\mathbf{1}\}$ and eigenvalue $1$ on
the complementary hyperplane, and $J = \kappa(I-P)$ has spectrum
$\{0\} \cup \{\kappa\}^{S-1}$. Since $\kappa \ge 0$, $\rho(J) = \kappa$.
\end{IEEEproof}

Three remarks. The zero eigenvalue is not an accident of the algebra: it lies
along $\mathbf{1}$, the direction in which every cell is scaled by a common
factor, and the map is invariant to exactly that rescaling because $\mu$ enters
only through $\mu_s/\bar{\mu}$. The iteration therefore has no contraction to
perform on the overall level and contracts at rate $\kappa$ on every
disparity-carrying direction. Second, $\rho(J)$ does not depend on $w$, so no
distribution of recorded rates -- however skewed -- moves the stability
threshold. Third, $\rho(J) < 1 \iff \kappa < 1$, which is the same threshold at
which the amplification exponent $1/(1-\kappa)$ of Theorem~1 has its pole; the
pole and the loss of contraction are one phenomenon.

We verified Lemma~\ref{lem:spectral} numerically by forming $J$ at the converged
fixed point and computing $\max_i|\lambda_i(J)|$ for
$\kappa \in \{0.30, 0.50, 0.70, 0.85, 0.95\}$ with $S = 400$ and
$\lambda \sim \mathrm{Gamma}(2,2) + 0.3$. The computed radii were
$0.300000$, $0.500000$, $0.700000$, $0.850000$, $0.950000$, matching $\kappa$ to
machine precision in every case.

\subsection{Attribution corollary: full proof and sign preservation}
\label{sec:s1-attribution}

Theorem~1(ii) of the main text gives the recorded-to-latent disparity map. We
restate the inverse and prove the two properties that determine how it may
honestly be used.

\begin{proposition}[Attribution and sign preservation]
\label{prop:attribution}
Let $\kappa \in [0,1)$ and let $\Delta_{\lambda} = \lambda_1/\lambda_2 > 0$ be a
between-group latent ratio with recorded counterpart $\Delta_y = \mu_1/\mu_2$
under the fixed point \eqref{eq:s-map}. Then
\begin{enumerate}
  \item[(i)] $\Delta_{\lambda} = \Delta_y^{\,1-\kappa}$, and the map
  $\Delta_y \mapsto \Delta_{\lambda}$ is a strictly increasing bijection of
  $(0,\infty)$ onto itself;
  \item[(ii)] the share of recorded log-disparity attributable to the loop is
  exactly $\kappa$, that is
  $\bigl(\log\Delta_y - \log\Delta_{\lambda}\bigr)/\log\Delta_y = \kappa$ for
  every $\Delta_y \ne 1$;
  \item[(iii)] $\Delta_{\lambda} > 1 \iff \Delta_y > 1$ and
  $\Delta_{\lambda} = 1 \iff \Delta_y = 1$.
\end{enumerate}
\end{proposition}

\begin{IEEEproof}
From the closed form derived in the main text,
$\mu_s = \lambda_s^{1/(1-\kappa)}M^{-\kappa/(1-\kappa)}$, the common factor
$M^{-\kappa/(1-\kappa)}$ cancels in the ratio, giving
$\Delta_y = \Delta_{\lambda}^{1/(1-\kappa)}$. Since $1-\kappa \in (0,1]$, the map
$x \mapsto x^{1/(1-\kappa)}$ is a strictly increasing bijection of $(0,\infty)$,
and its inverse is $x \mapsto x^{1-\kappa}$, which is (i). Taking logarithms in
(i) gives $\log\Delta_{\lambda} = (1-\kappa)\log\Delta_y$, hence
$\log\Delta_y - \log\Delta_{\lambda} = \kappa\log\Delta_y$; dividing by
$\log\Delta_y \ne 0$ gives (ii). For (iii), $x^{1-\kappa} > 1 \iff x > 1$ for any
positive exponent $1-\kappa > 0$, and equality transfers likewise.
\end{IEEEproof}

Part (iii) is the part we most want a reader to hold onto, because it is the
guard against a misuse the corollary otherwise invites. The loop inflates the
\emph{magnitude} of a disparity and can inflate it without bound as
$\kappa \to 1$, since $1-\kappa \to 0$ drives $\Delta_{\lambda} \to 1$ for any
fixed $\Delta_y$. It cannot create a disparity from nothing: a recorded gap
implies a real gap of the same direction, always. Any reading of
Proposition~\ref{prop:attribution} that concludes ``the observed disparity is an
artifact'' is a misreading; the correct conclusion is ``the observed disparity
overstates a real one by a factor the gain determines''.

Table~\ref{tab:s-attribution} gives the arithmetic. Each row was checked against
the numerically converged fixed point; agreement was within $10^{-12}$ relative
error at $\kappa \in \{0.3, 0.5, 0.6, 0.85\}$.

\begin{table}[!t]
  \centering
  \caption{Attribution map $\Delta_{\lambda} = \Delta_y^{1-\kappa}$. A recorded
  ratio on the left corresponds to the true ratio on the right at each gain. The
  final column is the share of recorded log-disparity attributable to feedback,
  which Proposition~\ref{prop:attribution}(ii) shows equals $\kappa$ exactly and
  is therefore independent of $\Delta_y$.}
  \label{tab:s-attribution}
  \begin{tabular}{r rrrr c}
    \toprule
    & \multicolumn{4}{c}{\textbf{true ratio $\Delta_{\lambda}$ at gain $\kappa$}}
    & \\
    \cmidrule(lr){2-5}
    \textbf{recorded $\Delta_y$} & $\kappa=0.3$ & $\kappa=0.5$ & $\kappa=0.6$
    & $\kappa=0.85$ & \textbf{loop share} \\
    \midrule
    $2{:}1$ & 1.63 & 1.41 & 1.32 & 1.11 & $\kappa$ \\
    $3{:}1$ & 2.16 & 1.73 & 1.55 & 1.18 & $\kappa$ \\
    $5{:}1$ & 3.09 & 2.24 & 1.90 & 1.27 & $\kappa$ \\
    $10{:}1$ & 5.01 & 3.16 & 2.51 & 1.41 & $\kappa$ \\
    \bottomrule
  \end{tabular}
\end{table}

\subsection{Why integer conformity scores make coverage conservative}
\label{sec:s1-lattice}

The main text reports that coverage converges to the nominal level from above and
attributes the residual excess to the integer lattice. That attribution is worth
making precise, because it explains a systematic feature of every coverage number
in the paper and it is the reason the randomized-PIT variant exists.

Two distinct discretizations are in play and they must not be conflated. First,
the target is a count: $y \in \mathbb{N}_0$. For any distribution $F$ on
$\mathbb{N}_0$ and any $p$, define the quantile
$Q_p = \min\{k \in \mathbb{N}_0 : F(k) \ge p\}$. Then
\begin{align}
\Pr\!\left(Q_{\alpha/2} \le y \le Q_{1-\alpha/2}\right)
&= F(Q_{1-\alpha/2}) - F(Q_{\alpha/2}-1) \notag\\
&\ge (1-\alpha/2) - \alpha/2 = 1-\alpha,
\label{eq:s-lattice1}
\end{align}
where the first inequality uses $F(Q_{1-\alpha/2}) \ge 1-\alpha/2$ by definition
of the quantile, and the second uses $F(Q_{\alpha/2}-1) < \alpha/2$, which holds
because $Q_{\alpha/2}$ is the \emph{smallest} integer whose CDF reaches
$\alpha/2$. The inequality in \eqref{eq:s-lattice1} is strict unless $F$ happens
to hit $\alpha/2$ and $1-\alpha/2$ exactly at lattice points, which for a Poisson
or ZINB law occurs on a measure-zero set of rates. So a discrete quantile
interval over-covers, and the excess is the probability mass of the two boundary
atoms.

Second, the CQR score inherits the lattice. With $q_{\alpha/2}, q_{1-\alpha/2}$
integer-valued and $y_i$ integer, the score
$V_i = \max(q_{\alpha/2}(x_i) - y_i,\, y_i - q_{1-\alpha/2}(x_i))$ takes values in
$\mathbb{Z}$. The split-conformal threshold is an order statistic of
$\{V_i\}$, hence also an integer, so the achievable coverage levels form a
discrete set: as the threshold steps from $t$ to $t+1$ the interval widens by two
counts and coverage jumps. There is in general \emph{no} integer threshold
attaining $1-\alpha$ exactly, and the split-conformal rule selects the smallest
integer whose empirical coverage is at least $1-\alpha$. Combining the two
effects, the finite-sample guarantee
\begin{equation}
\label{eq:s-lattice2}
1-\alpha \;\le\;
\Pr\!\left(y_{n+1} \in \widehat{C}_{n+1}\right)
\;\le\; 1-\alpha + \frac{1}{n+1}
\end{equation}
holds as usual, but the realized coverage in \eqref{eq:s-lattice2} sits near the
\emph{upper} end rather than at the nominal value, and the gap does not vanish as
$n \to \infty$: the $1/(n+1)$ term does, while the lattice term does not.

This is visible in our results in two independent ways. Nine of the ten
calibrators in the main text over-cover relative to the $90\%$ target, by between
$0.02$ and $4.05$ percentage points; the tenth, rolling adaptive ECRC, under-covers
in both cities ($0.8930$, $0.8918$) and is rejected by the selection policy for
precisely that reason, despite having the best disparity and the narrowest
intervals. And on NYC, four calibrators -- split CP, recency-weighted CP,
Mondrian and equalized coverage -- return numerically identical coverage, width
and disparity ($0.9326$, $18.57$, $0.0201$), because their differing quantile
rules select the same integer. Methods that differ only in \emph{how} they choose
a quantile level cannot differ in output when the level maps to the same lattice
point. The randomized-PIT variant restores a continuous scale by construction and
is the correct object on which to check exactness; the integer intervals it
delivers still over-cover, and the main text reports both numbers.

\subsection{Sensitivity to a mismeasured feedback gain}
\label{sec:s1-sensitivity}

Theorem~2(iii) of the main text states the log-error of the deflation under an
estimated gain. The derivation is short and we give it in full, together with the
step that turns an unbounded error into a bounded one on the retained set.

The deflation at gain $\hat{\kappa}$ is
$\hat{\lambda}_s = \hat{\mu}_s / (\hat{\mu}_s/M)^{\hat\kappa}$, so
\begin{equation}
\label{eq:s-sens1}
\log\hat{\lambda}_s = \log\hat{\mu}_s - \hat\kappa\,\log(\hat{\mu}_s/M).
\end{equation}
By Theorem~2(i), at the true gain the same expression \eqref{eq:s-sens1} returns
the truth:
$\log\lambda_s = \log\mu_s - \kappa\log(\mu_s/M)$. Subtracting, and treating the
recorded rate as given so that $\hat\mu_s = \mu_s$,
\begin{equation}
\label{eq:s-sens2}
\log\hat{\lambda}_s - \log\lambda_s
= (\kappa - \hat\kappa)\,\log\!\left(\mu_s/M\right).
\end{equation}
The right-hand side is a product of an error in the gain and a cell-specific
leverage $\log(\mu_s/M)$, and the leverage is unbounded: a cell recorded far from
the panel mean amplifies any error in $\kappa$ without limit. That is the precise
sense in which the correction is not uniformly safe, and it is what the abstention
rule exists to bound.

On the retained set
$\mathcal{R} = \{s : \hat\kappa < \kappa_{\mathrm{run}},\;
m_s \in [\bar{m}^{-1}, \bar{m}]\}$ the multiplier constraint
$m_s = (\mu_s/M)^{\hat\kappa} \in [\bar{m}^{-1},\bar{m}]$ is equivalent, on taking
logarithms, to $|\hat\kappa \log(\mu_s/M)| \le \log\bar{m}$, hence
\begin{equation}
\label{eq:s-sens3}
\left|\log(\mu_s/M)\right| \;\le\; \frac{\log\bar{m}}{\hat\kappa},
\end{equation}
and substituting \eqref{eq:s-sens3} into \eqref{eq:s-sens2},
\begin{equation}
\label{eq:s-sens4}
\left|\log\hat{\lambda}_s - \log\lambda_s\right|
\;\le\; \left|\kappa - \hat\kappa\right| \frac{\log\bar{m}}{\hat\kappa}
\qquad \text{for all } s \in \mathcal{R}.
\end{equation}
With the implementation values $\bar{m} = 5$ and, say, $\hat\kappa = 0.85$, the
leverage factor is $\log 5 / 0.85 = 1.893$, so a gain error of $0.01$ costs at
most a $1.9\%$ relative error in the deflated rate. The bound is monotone
decreasing in $\hat\kappa$, which is at first counter-intuitive: a larger gain
gives a \emph{tighter} bound per unit of gain error. The resolution is that the
constraint $m_s \in [\bar{m}^{-1},\bar{m}]$ is itself harder to satisfy at large
$\hat\kappa$, so the retained set shrinks. The cost does not disappear; it moves
from interval width into abstention rate, which is why the main text reports the
retained fraction beside every corrected coverage number.

Note finally what \eqref{eq:s-sens4} does not cover. It bounds the error from a
wrong \emph{gain} within a correctly specified power-law recording map. An error
in the functional \emph{form} of that map is a different failure and is treated by
the $\Gamma$ model of Section~\ref{sec:s2}.

% =====================================================================
\section{Extended $\Gamma$-Sensitivity Frontier}
\label{sec:s2}
% =====================================================================

The main text reports five points on the sensitivity frontier. Here we give the
full sweep and state exactly what the construction does and does not buy.

The sensitivity model is Rosenbaum's, transplanted from treatment assignment to a
recording mechanism. Suppose the true multiplier lies within a factor
$\Gamma \ge 1$ of the assumed one at every cell,
\begin{equation}
\label{eq:s-gamma}
m^{\mathrm{true}}_s / m_s \in [\Gamma^{-1}, \Gamma]
\qquad \text{for all } s .
\end{equation}
Because $\lambda_s = \mu_s/m^{\mathrm{true}}_s$ and
$\hat{\lambda}_s = \mu_s/m_s$, \eqref{eq:s-gamma} transfers directly to the rate,
$\lambda_s/\hat{\lambda}_s \in [\Gamma^{-1},\Gamma]$, and the inflated interval
$[\,Q^{\mathrm{Pois}}_{\alpha/2}(\hat{\lambda}_s/\Gamma),\;
Q^{\mathrm{Pois}}_{1-\alpha/2}(\hat{\lambda}_s\Gamma)\,]$ contains the
un-inflated interval at the true rate for every admissible
$m^{\mathrm{true}}$. Coverage therefore holds simultaneously over the whole band,
and $\Gamma = 1$ recovers the plain interval.

Table~\ref{tab:s-gamma} reports the frontier. Each row is the mean over $8$
trials of $6000$ cells at true gain $\kappa = 0.6$ and target $1-\alpha = 0.90$,
in worlds whose true multiplier deviates from the assumed one by a per-cell factor
drawn log-uniformly on $[\Gamma^{-1},\Gamma]$. Reproduced by
\texttt{scripts/misspecification\_sensitivity.py}; raw values in
\texttt{outputs/gamma\_sensitivity\_frontier.json}.

\begin{table}[!t]
  \centering
  \caption{Sensitivity frontier at $\kappa = 0.6$, target $0.90$. ``Plain'' is
  the un-inflated corrected interval evaluated in a misspecified world;
  ``inflated'' applies the matching $\Gamma$ band. Width is relative to
  $\Gamma = 1$. The plain interval last meets target at $\Gamma = 1.4$
  ($0.9057$), fails from $\Gamma = 1.6$, and reaches $0.6762$ at $\Gamma = 3$;
  the inflated interval never falls below target.}
  \label{tab:s-gamma}
  \begin{tabular}{r rr r l}
    \toprule
    $\Gamma$ & \textbf{plain coverage} & \textbf{inflated coverage}
    & \textbf{width $\times$} & \textbf{plain vs. target} \\
    \midrule
    1.0 & 0.9420 & 0.9420 & 1.00 & meets \\
    1.2 & 0.9328 & 0.9775 & 1.27 & meets \\
    1.4 & 0.9057 & 0.9872 & 1.51 & meets \\
    1.6 & 0.8802 & 0.9922 & 1.72 & \textbf{fails} \\
    1.8 & 0.8457 & 0.9939 & 1.94 & \textbf{fails} \\
    2.0 & 0.8091 & 0.9945 & 2.14 & \textbf{fails} \\
    2.5 & 0.7320 & 0.9965 & 2.61 & \textbf{fails} \\
    3.0 & 0.6762 & 0.9975 & 3.07 & \textbf{fails} \\
    \bottomrule
  \end{tabular}
\end{table}

Three readings, one of them unfavourable to us. The frontier is the honest object
to report because it converts an unfalsifiable modelling assumption into a
priced one: a practitioner who believes the recording response could be off by up
to $60\%$ per cell can buy validity at $1.72\times$ the width, and one who
believes it could be off threefold pays $3.07\times$. Second, the breakdown point
of the un-inflated interval is informative in itself. Coverage survives
$\Gamma = 1.4$ and fails by $\Gamma = 1.6$, so the plain correction tolerates
roughly $40\%$ multiplicative misspecification before it stops being honest --
which is a modest tolerance and should be stated as such rather than buried.

Third, and this is the unfavourable part, the inflated intervals over-cover
substantially: $0.9775$ at $\Gamma = 1.2$ rising to $0.9975$ at $\Gamma = 3$
against a $0.90$ target. The $\Gamma$ band is a worst-case construction and pays
for its guarantee in exactly the currency conformal prediction is meant to
economize. A reader who needs sharp intervals under suspected misspecification
would be better served by a model of the recording response than by our envelope
around the absence of one. We report the envelope because it is available without
further assumptions, not because it is efficient.

% =====================================================================
\section{Baseline Architectures and Training Protocol}
\label{sec:s3}
% =====================================================================

Every specification in this section was read from
\texttt{scripts/deep\_baselines.py} and \texttt{scripts/baselines.py} at the
commit that produced the reported results, not reconstructed from memory. Where a
value is a constructor default that the calling code does not override, we say so.

\subsection{Deep spatiotemporal baselines}

All four consume the same inputs as CIVIC-SAFE: a window of $W = 52$ weeks of
counts over $C = 3$ categories concatenated with the static covariate vector, and
all four emit ZINB or NB parameters so that CRPS is computed identically across
methods. Table~\ref{tab:s-baselines} collects the configurations.

\begin{table}[!t]
  \centering
  \caption{Deep baseline configurations as instantiated. All are trained with the
  protocol of Section~\ref{sec:s3-protocol}.}
  \label{tab:s-baselines}
  \small
  \begin{tabular}{l l l}
    \toprule
    \textbf{Model} & \textbf{Configuration} & \textbf{Output / loss} \\
    \midrule
    LSTM-NB & 2 stacked LSTM layers, hidden $64$,
              dropout $0.1$, LayerNorm on the & NB $(\mu, r)$ via two
              MLP heads; \\
            & final hidden state, final timestep read out
              & negative binomial NLL \\
    \addlinespace
    TFT-ZINB & variable-selection network with gated residual networks per
               & ZINB $(\pi,\mu,r)$; \\
             & variable group, $d_{\mathrm{model}} = 64$, $4$ attention heads,
               $2$ blocks, & CRPS \\
             & dropout $0.1$ & \\
    \addlinespace
    GraphWaveNet & $32$ channels, $4$ WaveNet blocks with dilations
                   $1,2,4,8$, & ZINB $(\pi,\mu,r)$; \\
                 & kernel size $2$, adaptive adjacency from two learned
                   & CRPS \\
                 & node embeddings of dimension $16$, dropout $0.1$ & \\
    \addlinespace
    STZINB-GNN & graph convolution over the static adjacency followed by a
                 & ZINB $(\pi,\mu,r)$; \\
               & temporal GRU, hidden $64$ & CRPS \\
    \bottomrule
  \end{tabular}
\end{table}

Two details deserve stating rather than leaving to be inferred. GraphWaveNet's
receptive field follows from its dilation schedule: with kernel size $k = 2$ and
dilations $d_i = 2^i$ for $i = 0,\dots,3$, the field is
$1 + \sum_i (k-1)d_i = 1 + (1+2+4+8) = 16$ weeks. That is less than the $52$-week
input window, so the model cannot in principle see the annual cycle -- a genuine
architectural disadvantage relative to the seasonal-naive baseline, and one worth
noting when reading its CRPS. And the loss is not uniform across baselines:
LSTM-NB is trained on negative binomial NLL while the other three are trained
directly on CRPS. Since CRPS is the reported metric, LSTM-NB is evaluated on a
criterion it was not trained for. We did not equalize this, and a referee is
entitled to regard LSTM-NB's numbers as correspondingly pessimistic.

\subsection{Classical baselines}

\begin{itemize}
  \item \textbf{Rolling historical average.} The mean of the $52$-week input
  window, recomputed at every test week, so the forecast updates as new weeks
  arrive. This is the baseline the evaluation code marks as binding, and the one
  the main text reports skill against.
  \item \textbf{Frozen historical average.} A single mean over the training
  period, held fixed across the whole test year. Reported only to show how much
  larger the apparent skill becomes against a baseline no deployment would use.
  \item \textbf{Seasonal naive.} $\hat{y}(t) = y(t-52)$, the count from the same
  week one year earlier, read off the first slot of the $52$-week window.
  \item \textbf{Lag-1 persistence.} $\hat{y}(t) = y(t-1)$, the last slot of the
  window.
  \item \textbf{STARIMA.} An AR(1) model with a spatial-lag regressor, fitted by
  ordinary least squares separately per spatial unit and crime category.
  \item \textbf{Plain ZINB.} \texttt{statsmodels}
  \texttt{ZeroInflatedNegativeBinomialP} with the covariate matrix in both the
  count and inflation components, fitted by BFGS with \texttt{maxiter=30}. This
  is an unregularized maximum-likelihood fit with no floor on the dispersion
  parameter, and it diverged on NYC (CRPS $924.10$); the main text excludes it
  from ranking rather than claiming credit for the gap.
  \item \textbf{XGBoost.} Gradient-boosted trees on the flattened window,
  producing point forecasts.
\end{itemize}

All seven classical baselines emit a point forecast, and CRPS is a
distributional score, so the two must be reconciled. Our evaluation dresses each
point forecast in a Poisson predictive distribution at that point --
$\pi = 0$, $\mu = \hat{y}$, $r = 1000$ so the negative binomial collapses to a
Poisson -- and scores it with the same 	exttt{crps\_zinb} routine used for every
other method. This is a deliberately generous convention and it is not neutral: a
Poisson at the right mean has appropriate spread for count data, so it scores far
better than the degenerate point mass a naive reading would assign. The effect is
large. XGBoost on Chicago has MAE $3.9688$ but CRPS $2.9157$; the gap is entirely
the Poisson dressing.

This matters for the headline comparison rather than being a technicality. When we
report that the rolling historical average beats all four deep spatiotemporal
baselines, the rolling average is being scored as a Poisson forecaster at the
trailing mean, not as a bare point predictor. That is the right comparison to make
-- it is what a competent deployment of a moving average would actually issue --
but a reader should know that the baseline was given a well-specified
distributional form for free, and that a less charitable convention would have
flattered our method instead.

\subsection{Training protocol and environment}
\label{sec:s3-protocol}

Optimizer AdamW, learning rate $10^{-3}$, weight decay $10^{-5}$, cosine
annealing over the epoch budget, early stopping on validation CRPS with patience
$10$. Deep baselines are trained for three seeds, $\{42, 137, 256\}$, and reported
as mean $\pm$ standard deviation. CIVIC-SAFE uses five seeds,
$\{42, 137, 256, 512, 1024\}$, combined by category-conditioned
entropy-regularized EMOS with a fallback to uniform weighting when the learned
weights fail to beat uniform on an internal calibration holdout.

The seed budget is uneven between our method and the baselines, and this matters
for the comparison: five seeds of ensembling is doing substantial work for us, as
the main text states -- a single CIVIC-SAFE seed averages CRPS $3.3622$ and loses
to TFT-ZINB's $2.9456$ on Chicago. A fully symmetric study would ensemble the
baselines to five seeds as well. We did not, and the honest reading of the
head-to-head result is that it compares an ensembled method against
un-ensembled ones. Our configuration also notes that seven to eight seeds would
be needed for paired $t$-tests at power $0.80$ with effect size $d = 0.8$; we ran
five.

% =====================================================================
\section{Dataset, Covariate and Graph Audits}
\label{sec:s4}
% =====================================================================

\subsection{Demographic covariates}

The demographic panel is built from American Community Survey five-year estimates
and contains, per spatial unit, exactly these fields:
\texttt{total\_population}, \texttt{median\_household\_income},
\texttt{poverty\_rate}, \texttt{unemployment\_rate}, \texttt{pct\_black},
\texttt{pct\_hispanic}, \texttt{pct\_renter\_occupied} and
\texttt{population\_density}. Seven are used as covariates with total population
serving as the normalizer for density and rate construction.

We flag one absence explicitly because it is easy to assume otherwise:
\textbf{there is no educational-attainment variable} in the panel. Any reading of
our fairness audit that supposes education is controlled for is mistaken. The
demographic stratification used throughout the fairness analysis is a quartile
split on \texttt{pct\_black}, so the ``higher-minority stratum'' of the main text
means units above the median on that single covariate, not a composite index.

\subsection{Spatial graph construction and connectivity}

Two adjacency structures are built per city. Queen contiguity joins polygons that
share any boundary point, computed by spatial join on geometries snapped to a
$10^{-4}$-degree precision grid to defuse floating-point boundary comparisons,
with an \texttt{intersects} fallback if \texttt{touches} yields fewer edges than
nodes. The second graph joins each unit to its $k = 8$ nearest neighbours by
centroid distance in a projected meter-based coordinate reference system
(EPSG:26971 for Chicago, EPSG:32118 for New York).

As specified in the main text, the second graph is \textbf{geographic}, not demographic:
it connects spatially proximate units that do not happen to share a boundary, and it does
not consult the covariate vector at all. Describing it as a nearest-neighbour graph in
demographic feature space would misstate the construction and also misstate its motivation --
what it supplies is proximity beyond adjacency, not similarity of composition.

Table~\ref{tab:s-graphs} reports the realized connectivity. The measured
statistics justify the dual-graph design more concretely than the motivation given
in the main text: under queen contiguity alone, \textbf{one New York precinct is
isolated}, with degree zero, and would receive no spatial message passing
whatever. The $k$-nearest-neighbour graph raises its degree to at least two.
Chicago has no isolated unit but its minimum queen degree is $1$, so the same
mechanism is operating less dramatically there.

\begin{table}[!t]
  \centering
  \caption{Realized graph connectivity, computed from the released shapefiles by
  \texttt{build\_adjacency\_from\_geodataframe} with $k = 8$. Degrees are
  in-degrees. The isolated NYC precinct under queen contiguity is the concrete
  reason the second graph is not redundant.}
  \label{tab:s-graphs}
  \begin{tabular}{l l r r r r r}
    \toprule
    \textbf{City} & \textbf{Graph} & \textbf{Nodes} & \textbf{Edges}
    & \textbf{Mean deg.} & \textbf{Min} & \textbf{Max} \\
    \midrule
    Chicago & queen & 77 & 394 & 5.12 & 1 & 9 \\
    Chicago & $k$-NN ($k{=}8$) & 77 & 616 & 8.00 & 1 & 13 \\
    \addlinespace
    New York & queen & 78 & 240 & 3.08 & \textbf{0} & 6 \\
    New York & $k$-NN ($k{=}8$) & 78 & 624 & 8.00 & 2 & 13 \\
    \bottomrule
  \end{tabular}
\end{table}

The $k$-NN mean in-degree is exactly $8.00$ in both cities because every node
contributes exactly eight out-edges; the in-degree spread of $1$ to $13$ reflects
the asymmetry of nearest-neighbour relations, and the graph is used as a directed
edge set without symmetrization.

% =====================================================================
\section{Additional Diagnostics}
\label{sec:s5}
% =====================================================================

\subsection{Group-stratified coverage}

Table~\ref{tab:s-groups} reports coverage and mean width within each demographic
quartile for the calibrator selected in each city. The disparity figures quoted in
the main text, $0.0238$ and $0.0286$, are the ranges of the coverage column.

\begin{table}[!t]
  \centering
  \caption{Coverage and width by demographic quartile for the selected
  calibrator. Group $0$ is the lowest \texttt{pct\_black} quartile. Abstention is
  zero in every cell, so $n_{\mathrm{issued}} = n$ throughout.}
  \label{tab:s-groups}
  \begin{tabular}{l l r r r}
    \toprule
    \textbf{City} & \textbf{Group} & $n$ & \textbf{Coverage}
    & \textbf{Mean width} \\
    \midrule
    Chicago & 0 & 3180 & 0.8994 & 17.65 \\
    (equalized & 1 & 3021 & 0.9037 & 13.36 \\
    coverage) & 2 & 3021 & 0.9043 & 12.71 \\
            & 3 & 3021 & 0.9232 & 14.43 \\
    \addlinespace
    New York & 0 & 3180 & 0.8858 & 16.94 \\
    (variance-scaled & 1 & 3021 & 0.9020 & 19.60 \\
    split CP) & 2 & 3180 & 0.9145 & 14.94 \\
            & 3 & 3021 & 0.8984 & 14.35 \\
    \bottomrule
  \end{tabular}
\end{table}

Two observations that the marginal numbers conceal. Coverage is not uniformly at
or above nominal within groups: Chicago group $0$ sits at $0.8994$ and New York
groups $0$ and $3$ at $0.8858$ and $0.8984$, all below the $0.90$ target. The
marginal guarantee is marginal, and a group-conditional guarantee is what ECRC and
the Mondrian variants target at the cost of width. And width varies more across
groups than coverage does -- Chicago spans $12.71$ to $17.65$, a factor of
$1.39$ -- so equal coverage is being purchased with unequal sharpness. Whether
that trade is acceptable is a policy question our constraint does not resolve; we
report it because a coverage-only audit would hide it.

\subsection{Per-category decomposition}

Table~\ref{tab:s-categories} gives the per-category CRPS for the model and the
rolling historical average, with the resulting skill scores.

\begin{table}[!t]
  \centering
  \caption{Per-category CRPS and skill score against the rolling historical
  average, Chicago. Categories differ by more than an order of magnitude in
  scale, so the aggregate CRPS is dominated by property crime.}
  \label{tab:s-categories}
  \begin{tabular}{l r r r}
    \toprule
    \textbf{Category} & \textbf{CIVIC-SAFE CRPS} & \textbf{Rolling HA CRPS}
    & \textbf{CRPSS} \\
    \midrule
    Violent & 3.1324 & 3.3220 & $+0.0571$ \\
    Property & 4.7967 & 4.9166 & $+0.0244$ \\
    Drug & 0.5510 & 0.5579 & $+0.0124$ \\
    \bottomrule
  \end{tabular}
\end{table}

The decomposition is more informative than the aggregate. Skill is concentrated in
violent crime, at $+5.71\%$, and is close to negligible on drug crime at
$+1.24\%$ -- which is the sparsest category, at $50.89\%$ zero cells, and
therefore the one where the zero-inflated component of the likelihood is doing the
most work and the least room exists to beat a trailing mean. The aggregate
$+3.60\%$ is a scale-weighted average dominated by property crime, the category on
which our margin is second smallest. A reader interested in violent-crime
forecasting specifically should take $+5.71\%$; a reader interested in the
aggregate should take $+3.60\%$; neither should take the largest number in the
paper.

\subsection{A diagnostic we cannot supply}

A group-stratified probability integral transform -- PIT uniformity tested
separately within each demographic quartile -- would sharpen the Chicago
misspecification finding considerably, by establishing whether the compressed
right tail is concentrated in particular strata or is uniform across them. We
cannot compute it from the released artifacts. The randomized PIT requires the
per-cell ZINB parameters $(\pi_s, \mu_s, r_s)$, and the persisted prediction
file carries only observed counts, point predictions, the rolling-average
reference, the conformal upper bound and the group label. Recomputing it requires
re-running inference from the checkpoints and persisting the parameter tensors,
a one-line change to the evaluation script that we have not made.

We state this rather than substitute the pooled PIT and let the section heading
imply more than was done. The pooled result stands: Chicago
$\chi^2 = 241.84$ on nine degrees of freedom, $p = 5.25\times10^{-47}$; New York
$\chi^2 = 10.75$, $p = 0.293$. The spatial residual map in the main text is the
closest available proxy for the stratified question, and it shows the shortfall
concentrated on the West Side rather than distributed evenly, with $65$ of $77$
areas under-predicted and a correlation of $-0.66$ between an area's mean count
and its mean residual.

\end{document}
